\scanpage{274}
\mspage{- 80 -}\hfill IX

6.4. Notons que lorsque $\ell\ne p$, nous avons mis en évidence dans 4.3 a) une fonction à valeurs entières remarquable sur le groupe d'inertie $I$ de $S$, d'ailleurs indépendante du choix de $\ell$, et dont la connaissance, dans le cas d'un corps résiduel parfait, suffit à reconstruire l'exposant conducteur de $A$ (4.5), en prenant le produit scalaire, sur un quotient assez grand $\Gamma=I/I'$ de $I$, de la fonction précédente et du conducteur d'Artin sur $\Gamma$ (cf. (4.6.2)). Il peut être intéressant d'observer que cette fonction à valeurs entières sur $I$ peut se réconstituer également par la connaissance de $T_\ell(A_K)$ pour $\ell=p$ - du moins (pour le moment) en inégales caractéristiques. Lorsque le corps résiduel $k$ est fini, une observation analogue s'applique à la fonction à valeurs dans $\mathbb Q$, définie sur le sous-groupe $\pi^!$ de $\pi$ formé des éléments dont l'image dans $\pi_{\unsure{o|0}}$ est de la forme $f_q^n$ ($n\in\mathbb Z$), considérée dans 4.3 b). Indiquons le principe de la vérification de ces assertions. On utilise la filtration (6.3.1), et on est réduit à attacher une fonction convenable (sur $I$ resp. sur $\pi^!$) à un $p$-pro-groupe de Barsoti-Tate $Y$ sur $K$ qui a "essentiellement bonne réduction sur $S$ ", au sens précisé dans 6.5. ($Y$ sera un des trois quotients provenant de la filtration (6.3.1).) Dans le cas de 4.3 a), on se \unsure{ramène|remène} au cas d'un corps résiduel parfait (ou même algébriquement clos, si on veut) par extension résiduelle convenable dans $S$ (EGA $0_{\mathrm{III}}$ 10.3.1). Prenons alors une extension galoisienne finie $K'$ de $K$ telle que $Y_{K'}$ ait bonne réduction sur $S'$, d'où un groupe de Barsotti-Tate $Y_{S'}$ sur $S'$ et une opération de $G$ sur $Y_{S'}$ compatible avec son action sur $S'$. Si $k'$ est le corps résiduel de $S'$, $G$

% Reverse impression from scan 273: the backward R in the upper margin and the marks across « on se ramène au cas ».
% archive sheet number: 273
% stamp: 398 (lower centre label; the pencilled 398 above is also visible).
